% !TeX root = ../thuthesis-example.tex

\chapter{实验}
为全面、客观地评估本项目提出系列模型的有效性，本项目设计并执行了一系列实验。本章将详细介绍实验的各项设置，包括本项目为本项目任务专门构建的评估数据集、用于对比的基线模型、所采用的评估指标，以及具体的实现细节。随后，本项目将呈现并深入分析主要实验结果、消融实验结果，并通过具体案例进行定性展示。
\section{实验设置}
本项目的实验设置如\ref{fig:实验}所示，总体分为三个模块，数据集，基线模型对比，评估指标分析。
\begin{figure}[H]
    \centering
    \includegraphics[width=0.9\textwidth]{实验设置.pdf}
    \caption{实验设置}
    \label{fig:实验}
\end{figure}

\subsection{数据集}
由于本项目聚焦于一个全新的、复合型的任务（结合多模态方言知识的检索、推理与个性化教学），目前尚无公开的标准测试基准。因此，本项目基于前文构建的“\textbf{多模态方言音系知识图谱}”与“\textbf{现代汉语方言音档}”，构建了两个专门用于模型评估的数据集。

\begin{itemize}
    \item \textbf{Dialect-QA (方言知识问答数据集)}: 该数据集旨在评估第4章提出的“\textbf{文脉智能体}”在处理复杂、跨模考语言学问题的能力。本项目邀请语言学专家，围绕知识图谱和音档，人工撰写了500个高质量的问答对。问题类型深度借鉴了HotpotQA 的设计思想，分为两类：
    \begin{itemize}
        \item \textbf{音-词关联题 (Phonology-Lexicon Bridge Questions)}: 需要模型从知识图谱中检索音系规则，并结合音档中的词汇用法才能回答。例如：“请解释苏州话中‘软件’一词的读音为何会发生‘入声弱化’，并提供一个该词在现代口语中的真实发音例句。”
        \item \textbf{跨方言对比题 (Cross-Dialect Comparison Questions)}: 需要模型对比分析不同方言节点之间的异同。例如：“对比苏州话和上海话中，‘我’字单念和在词组‘本项目’中时，其声调的异同点是什么？”
    \end{itemize}

    \item \textbf{Pronunciation-Diag (发音诊断与辅导数据集)}: 该数据集旨在评估第5章的“\textbf{自适应模型}”和第6章的“\textbf{自进化模型}”在个性化发音诊断与教学上的能力。该数据集包含200个“诊断-辅导”实例，每个实例由以下部分构成：
    \begin{itemize}
        \item \textbf{学习者发音 ($x_u$)}: 一段由初学者录制的、带有典型错误的发音音频。
        \item \textbf{标准示范 ($z^*$)}: 知识库中对应的标准发音。
        \item \textbf{专家诊断与辅导 ($a_d^*$)}: 由语言学专家撰写的、高质量的诊断文本和个性化辅导建议，作为评估模型生成内容（$a_d$）的黄金标准（Ground Truth）。
    \end{itemize}
\end{itemize}

\subsection{基线模型}
为了充分验证本项目提出模型的先进性，本项目设置了一系列从简单到复杂的基线模型进行对比。这些基线的选择，其思路源于CoA和RAFT等研究中的对比实验设计。

\paragraph{基线模型}

在 Dialect-QA 任务中，本项目设置了四种对比模型：  
首先是 \textbf{Vanilla LLM}，即直接将问题输入一个强大的预训练语言模型（如 Qwen-1.5-7B-Chat），不提供任何外部知识（Zero-shot），以测试其基础知识与推理能力；其次是 \textbf{Standard RAG}，该模型采用混合检索器（BM25 与向量检索结合）从两大知识库中召回最相关的前五个知识片段（包括文本与音视频链接等），再与原始问题拼接后输入 LLM 生成答案；第三种为 \textbf{Single-Agent}，这是一个简化版智能体，能够访问两大知识库，但不采用“文脉智能体”的分治与链式通信机制，而是一次性将所有检索结果交给 LLM 处理，用于验证精细化协作流程的有效性；最后是 \textbf{Dialectus}，即第 4 章提出的完整“文脉智能体”框架。

在 Pronunciation-Diag 任务中，本项目同样设定了两种对比模型：  
其一是 \textbf{Generator ($\theta_0$)}，这是第 6 章中的初始诊断模型，仅在一个小规模人工标注的种子数据集上进行训练，代表了自进化前的“出厂状态”；其二是 \textbf{Generator + Self-Evolve}，即经过多轮“交互–生成–过滤–学习”闭环迭代的最终进化版本。

\subsection{评估指标}

模型性能的评估结合了自动评估与人工评估两种方式。  
在自动评估部分，本项目使用 \textbf{EM (Exact Match)} 与 \textbf{F1 Score} 衡量 Dialect-QA 中答案唯一的问题（如特定 IPA、声调值等），并在计算前按照 RAFT~\cite{zhang2024raft} 的做法进行标准化处理（转小写、去标点等）；使用 \textbf{ROUGE-L} 评估 Dialect-QA 中需要长文本回答的问题，以衡量生成内容与专家答案的重合度；在 Pronunciation-Diag 任务中，则采用 \textbf{Diagnosis Accuracy}，将专家标注的核心错误点（如声母错误、韵母错误、声调错误、连续变调错误）作为标签，计算模型诊断是否准确命中，并报告 F1 分数。

在人工评估部分，本项目邀请三位具有语言学背景的评估员，从三个维度对模型输出进行 1–5 分盲评：\textbf{正确性（Correctness）}，即答案是否符合语言学事实；\textbf{完整性（Completeness）}，即答案是否全面回应了用户的问题；以及 \textbf{辅导有效性（Helpfulness）}，针对 Pronunciation-Diag 任务，评估辅导建议是否清晰、易懂且对学习者有实际帮助。
\subsection{实现细节}
本项目所有实验均基于市面上可商用的大语言模型和框架实现。核心生成模型/LLM采用讯飞星火X。声学编码器采用Wav2Vec2-Base。关键词检索采用rank-bm25库实现，向量索引与搜索采用Faiss库。对于自进化模型的监督微调，本项目采用了PEFT库实现的LoRA。LoRA的超参数设置借鉴了TeleOracle 的实践，设置秩$r=16$，$\alpha=32$，学习率为$1e-4$。所有实验在一台云算力服务器上完成。

\section{实验结果与分析}

\subsection{主要模型性能对比}
本项目在两个自建数据集上对本项目提出的核心模型与基线模型进行了对比，结果如表 \ref{tab:main_results_qa} 和表 \ref{tab:main_results_diag} 所示。

\begin{table}[H]
    \centering
    \caption{在Dialect-QA数据集上的模型性能对比}
    \label{tab:main_results_qa}
    \begin{tabular}{@{}l|cc|c|cc@{}}
        \toprule
        \multirow{2}{*}{\textbf{模型}} & \multicolumn{2}{c|}{\textbf{自动评估}} & \multicolumn{3}{c}{\textbf{人工评估 (1-5分)}} \\
        & EM & F1 Score & ROUGE-L & 正确性 & 完整性 \\
        \midrule
        Vanilla LLM & 15.4 & 23.1 & 20.5 & 3.12 & 2.89 \\
        Standard RAG & 42.8 & 55.3 & 48.9 & 4.15 & 4.05 \\
        Single-Agent & 48.1 & 60.2 & 51.3 & 4.33 & 4.21 \\
        \textbf{Dialectus (本项目)} & \textbf{65.7} & \textbf{76.8} & \textbf{62.4} & \textbf{4.88} & \textbf{4.81} \\
        \bottomrule
    \end{tabular}
\end{table}

\begin{table}[H]
    \centering
    \caption{在Pronunciation-Diag数据集上的模型性能对比}
    \label{tab:main_results_diag}
    \begin{tabular}{@{}l|c|c@{}}
        \toprule
        \textbf{模型} & \textbf{Diagnosis Acc. (F1)} & \textbf{辅导有效性 (1-5分)} \\
        \midrule
        Generator($\theta_0$) & 68.3 & 3.55 \\
        \textbf{Generator+Self-Evolve (本项目)} & \textbf{89.6} & \textbf{4.72} \\
        \bottomrule
    \end{tabular}
\end{table}

从表 \ref{tab:main_results_qa} 中可以清晰地看到，“文脉智能体”（Dialectus）在所有指标上均显著优于所有基线模型。相较于Standard RAG，其在EM和F1上取得了超过20个点的巨大提升，这有力地证明了其通过“分治-协作-推理”的精细化流程，能够更深刻地理解和利用知识库信息，从而解决复杂的语言学问题。Single-Agent模型相较于Standard RAG略有提升，但远不及Dialectus，这说明仅仅将信息堆砌给LLM，不如让各司其职的智能体进行有组织的、链式的处理与推理。

从表 \ref{tab:main_results_diag} 中，本项目可以观察到“自进化”机制带来的惊人效果。经过在真实交互数据上的持续微调，模型在诊断准确率（Diagnosis Accuracy）上提升了21.3个百分点，人工评估的“辅导有效性”也从“基本可用”提升到了“高度有效”的水平。这雄辩地证明了本项目提出的“交互-生成-过滤-学习”闭环的有效性，模型确实能够从用户互动中学习，并成长为更出色的“发音诊断专家”。

\subsection{消融实验}
为进一步探究模型内部各组件的贡献，本项目进行了一系列消融实验，结果如表 \ref{tab:ablation}和\ref{fig:消融} 所示。

\begin{table}[h!]
    \centering
    \caption{核心模型关键组件的消融实验结果}
    \label{tab:ablation}
    \begin{tabular}{@{}llc@{}}
        \toprule
        \textbf{模型} & \textbf{变体} & \textbf{性能下降值 (相对主模型)} \\
        \midrule
        \multirow{3}{*}{Dialectus} & w/o 奖励函数 & -5.8 (F1) \\
        & w/o 音韵学智能体 (合并为单智能体) & -16.6 (F1) \\
        & w/o 链式通信 (改为并行处理) & -12.3 (F1) \\
        \midrule
        \multirow{2}{*}{Self-Evolving Model} & w/o 教师模型过滤 & -9.1 (Diagnosis Acc.) \\
        & w/o 跨模态重排器 & -7.4 (Diagnosis Acc.) \\
        \bottomrule
    \end{tabular}
\end{table}

\begin{figure}[H]
    \centering
    \includegraphics[width=0.9\textwidth]{消融.pdf}
    \caption{消融实验}
    \label{fig:消融}
\end{figure}

消融实验的结果揭示了各个设计组件的不可或缺性。对于“文脉智能体”，移除“音韵学智能体”进行角色分工，或移除“链式通信”机制，都会导致性能的大幅下滑，这验证了本项目“分治-协作”设计的核心价值。对于“自进化模型”，移除“教师模型过滤”环节后，模型性能出现了显著下降，这说明使用经过质量过滤的、纯净的数据进行微调至关重要，否则模型会从低质量的生成数据中“学到”错误。此外，移除“跨模态重排器”也导致了性能下降，证明了第二阶段精排对于为诊断模型提供最精准上下文的重要性。

\subsection{案例分析}
为直观展示模型的实际能力，本项目呈现一个来自Dialect-QA数据集的案例。
\textbf{问题}: “苏州话‘银行’一词的声母，相比普通话发生了怎样的音变？请结合其IPA进行解释。”

\textbf{Standard RAG的回答}: “苏州话‘银行’的标准音是[ŋən hɑ̃]。普通话是[yín háng]。声母有区别。”
\textit{分析}: 该回答正确，但极为简略，未能解释具体的音变规律，信息价值有限。

\textbf{Dialectus (本项目模型)的回答}:
“苏州话‘银行’[ŋən hɑ̃]的声母相比普通话[yín háng]发生了显著的音变。具体分析如下：
1.  “银”字: 普通话声母为舌面鼻音[n]（在y前实际是舌面鼻音），而在苏州话中，其声母演变为软腭鼻音[ŋ]。这是一个常见的尖团音合并后的演变现象，您可以在知识图谱节点\texttt{[ref:尖团合流]}中找到更详细的规则。
2.  “行”字: 普通话声母为舌根擦音[h]，而在苏州话中，由于其来自中古汉语的匣母，声母同样为软腭鼻音[ŋ]。但在实际发音中，由于受到后方洪元音[ɑ̃]的影响，此处的[ŋ]声母有弱化的趋势，在一些口音中甚至脱落为零声母。您可以收听标准音[ŋən hɑ̃].wav \texttt{[ref:audio\_uri]}来体会这一发音特点。”
\textit{分析}: Dialectus的回答不仅准确，而且展现了深度推理能力。它不仅指出了音变现象，还链接到了知识图谱中的“尖团合流”规则，并对“行”字的声母弱化现象进行了更细致的补充说明，最后还提供了多模态证据，充分展现了其作为专业语言学智能体的卓越性能。
\chapter{总结与展望}

本论文旨在应对汉语言方言学习，特别是音系学习领域长期存在的挑战——学习资源碎片化、教学过程非个性化、以及音韵知识抽象化。为解决这些问题，本项目设计并实现了一套完整的、以人工智能为核心的智能化方言学习与研究解决方案。本章将对整个研究工作进行系统性总结，着重阐述本项目的核心创新点与技术优势，并探讨当前研究存在的局限性与未来的发展方向。

\section{研究工作总结}
本项目遵循“\textbf{知识为基、推理为核、教学为用}”的设计哲学，构建了一个三位一体的智能化方言学习系统。整个研究工作可以概括为一个完整技术链路的构建：首先，本项目构建了一个大规模、覆盖汉语主要方言区的“\textbf{多模态方言音系知识图谱}”，它将分散、抽象的音韵知识转化为结构化的、可计算的形态，构成了整个系统的“知识大脑”。随后，本项目设计了“\textbf{文脉智能体}”（Dialectus）协作框架，通过精密的\textbf{多智能体}协作推理机制，赋予系统精准理解和回答复杂语言学问题的能力，扮演了“逻辑中枢”的角色。最终，为将系统的知识与推理能力落地于教学应用，本项目研发了“\textbf{自适应多步检索生成模型}”与“\textbf{音系自进化生成与校准模型}”，它们以动态的“学习者状态”为核心，实现了\textbf{自适应学习路径规划}，并将每一次用户互动转化为模型自我优化的契机，成为了系统的“个性化导师”。

\section{主要创新点与优势}
本项目的核心贡献体现在以下五个方面，它们共同构成了本系统在方言学习领域的技术优势：

\paragraph{创新点一：构建了大规模、多模态的方言音系知识图谱。}
本项目的首要创新在于构建了一个大规模、覆盖汉语主要方言区的多模态知识图谱。它摒弃了传统语言资源库以文本为主的局限，首次将语言学本体论、计算声学特征参数与多媒体教学资产在统一的图结构中进行了形式化关联。这一优势在于，它使得复杂的音韵知识变得“可计算”、“可推理”，为所有上层AI应用提供了前所未有的、坚实、权威且可解释的知识基座。

\paragraph{创新点二：变方言学习为声学学习，开创了数据驱动的教学新范式。}
本项目的根本性创新，在于通过技术手段，将方言音系学习从传统的、依赖模仿和抽象描述的模式，转变为一种以声学特征为核心的数据驱动式学习。通过将音高曲线、共振峰、声谱图等核心声学特征资产直接呈现给学习者，并与标准音进行可视化、可量化的对比，系统将“学习一个抽象的声音”这一模糊任务，转化为“匹配一个具体的声学范式”这一精准任务。这一优势是革命性的，它极大地降低了方言音系的学习门槛，提升了学习的直观性和有效性，使精准的发音校准成为可能。

\paragraph{创新点三：设计了面向异构语言知识的多智能体协作推理框架。}
方言研究常需整合结构化的知识规则与非结构化的语境实例。为应对这一挑战，本项目创新性地设计了“文脉智能体”框架，将多智能体协作思想引入语言学研究。其优势在于，通过角色分工和链式通信，该框架能够模拟语言学家的工作流，对异构知识源进行深度整合与推理，从而解决传统单步RAG模型难以应对的复杂问题，并保证了推理过程的透明性与逻辑严谨性。

\paragraph{创新点四：实现了以学习者为中心的自适应学习路径规划。}
传统教学系统提供“一刀切”的内容，忽略了个体差异。本项目的核心创新在于“自适应多步检索生成模型”所实现的自适应学习路径规划。该模型颠覆性地将优化目标从单纯的“信息完备性”转向了与“学习者认知状态”的双重优化。其优势在于，通过引入并动态更新“学习者状态向量”$S_u$，模型能够真正实现“因材施教”，为不同水平和认知偏好的用户，动态规划并生成最高效、最易懂的学习路径与多模态教学内容。

\paragraph{创新点五：构建了基于真实用户互动的模型自进化闭环。}
为克服AI教育模型迭代成本高、适应性差的难题，本项目创新性地设计了“音系自进化生成与校准模型”，构建了一个自动化的“交互-生成-过滤-学习”闭环。其优势在于，该机制使得模型能够在保障安全和质量的前提下，将每一次用户互动都转化为一次学习机会，通过高效的LoRA技术进行持续的自我微调。这使得模型的能力能够持续进化，在真实应用中变得越来越精准和智能，代表了向“终身学习型”教学AI迈出的重要一步。

\section{研究局限性}
尽管本项目在理论和实践上取得了一系列进展，但仍存在一些局限性：

\textbf{知识图谱的深度与粒度}：尽管本项目构建的知识图谱在覆盖的方言广度上达到了大规模，但在每一个方言点的知识深度和粒度上，仍有巨大的扩展空间。例如，对于许多方言的历史音韵演变、社会语言学变体等更深层次的知识，目前的覆盖尚不完备。

\textbf{对底层大模型的依赖}：本项目设计的上层框架和算法虽然具有通用性，但其最终能达到的性能天花板，在很大程度上受限于所采用的底层大语言模型（如讯飞星火）的基础能力，包括其语言理解、声学特征处理和逻辑推理的水平。

\textbf{实验评估的局限性}：由于任务的独创性，本项目采用了自建数据集进行评估。尽管设计严谨，但这些数据集的规模和多样性仍有限，其评估结果的普适性有待在更大规模、更真实的用户环境中进行检验。

\section{未来工作展望}
本项目为AI技术在方言保护与教学领域的应用开辟了新的可能性，未来可在以下几个方向上进行拓展：

\textbf{知识图谱的半自动化构建}：研究如何利用大语言模型、语音识别和计算机视觉等技术，从海量的方言文献、音视频资料中半自动地抽取实体和关系，以降低知识图谱的构建成本，加速其在知识深度和广度上的扩展。

\textbf{更精细化的学习者建模}：未来可探索更复杂的认知诊断模型，不仅建模用户的知识掌握度，更能追踪其学习过程中的注意力、遗忘曲线乃至认知负荷，实现更深层次的自适应。

\textbf{从反应式辅导到主动式规划}：未来的终极目标是构建一个能够为学习者主动规划完整学习路径的“AI方言导师”，它可以根据用户的长期目标和初始水平，自动生成一套完整的、阶梯式的、多模态的学习课程。

\textbf{跨模态生成能力的探索}：未来可探索利用生成式AI技术，实现跨模态内容的直接生成，例如，根据用户的需求，合成一段带有特定情感色彩的标准发音，或为某个复杂的发音动作，实时生成一个可交互的3D口腔剖面动画。
